\documentclass[11pt]{article}
\usepackage[margin=1in]{geometry}
\usepackage{amsmath,amssymb,amsthm}
\usepackage{booktabs}
\usepackage{array}
\usepackage{enumitem}

\setlength{\parindent}{0pt}
\setlength{\parskip}{0.55\baselineskip}
\setlist{nosep,leftmargin=*}

\newtheorem{proposition}{Proposition}
\newtheorem{lemma}{Lemma}
\newcommand{\mstatus}[1]{\textbf{[Mathematical status: #1]}}
\newcommand{\istatus}[1]{\textbf{[Implementation applicability: #1]}}
\newcommand{\nstatus}[1]{\textbf{[Scientific novelty: #1]}}
\newcommand{\Prb}{\mathbb{P}}
\newcommand{\Ex}{\mathbb{E}}
\newcommand{\HARD}{\mathcal{H}}

\title{BLADE-NCO Theory Research\\Stage 1 with Stage 1.5 Formal-Model Corrections}
\author{Research note}
\date{October 21, 2026}

\begin{document}
\maketitle

\begin{abstract}
This note formalizes structural reachability, inherited population memory, moving hard-instance regions, and first discovery measured in computational cost. It records three candidate claims, attempts a proof of a finite-window memory comparison, and tests broad memory/diversity claims against exact finite-state examples. The strongest surviving result is conditional: retained structural memory helps only when it supplies a verified per-cost hit-hazard advantage over the comparator. No unconditional BLADE-NCO advantage is established. This is a Stage 1 note, not a completed theory paper.

Stage 1.5 corrects the state space, time indices, probability space,
aggregation language, and implementation mapping without adding a theorem or
changing that conclusion. The finite examples remain counterexamples to
universal claims; they are not assumptions that the TSP instance space is
finite.
\end{abstract}

\section{Scope and status}
The scope is restricted to structural reachability, evolutionary memory, and expected first-hitting time under a non-stationary hardness landscape. The supplied BLADE-NCO context is FOGA-11 structural mutation, inherited parent populations, HGB40 hardness estimation, Elite-QD1000 (P=T), and BQ solver co-adaptation. Those names are treated as labels, not as complete algorithm specifications.

Statuses are reported on three separate axes. \mstatus{PROVED}, for example,
concerns the validity of a mathematical statement under its premises;
\istatus{VERIFIED, PARTIAL, OPEN, or NOT APPLICABLE} concerns whether those
premises and variables match the inspected S6-AE implementation; and
\nstatus{OPEN or NOT CLAIMED} concerns literature priority. A proof of a
conditional probability bound does not establish implementation applicability
or novelty. Compilation verifies only that the source typesets; it does not
verify mathematical correctness.

The Stage 1.5 implementation boundary is the verified HGB40 + Elite-QD1000
$(P=T)$ S6-AE path. Its frozen artifacts passed 40/40 strict checks. Historical
GP4 output is not used as the definition of the current method. Details that
could not be established from the inspected code and artifacts are marked
OPEN rather than inferred.

\section{Prior-theory boundary}
Related theory includes rigorous results in which populations track moving optima and in which diversity mechanisms have different effects, but those works analyze explicit dynamic-function classes and specified algorithm dynamics \cite{Dang2017,OlivetoZarges2015}. They do not establish compute-normalized discovery of solver-relative hard instances under solver co-adaptation. Thus a generic claim that ``memory/diversity helps in dynamic optimization'' is unsuitable as a novelty claim. The present abstraction differs in its target: discovering instances whose measured hardness changes as the solver co-adapts, under a compute budget and with a surrogate-based archive. Whether that narrower combination is novel remains \nstatus{OPEN}; this note is not a literature-priority search.

\section{Formal model and implementation boundary}

\subsection{Structural states and mutation}
Let $\mathcal{Z}$ be the set of encoded candidate instances and let
$\phi:\mathcal{Z}\to\mathcal{X}$ map an encoding to the solver-visible
structural phenotype. Encodings $z,z'$ are equivalent when
$\phi(z)=\phi(z')$. A structural state is a phenotype $x\in\mathcal{X}$,
for example a canonical graph/instance together with whichever structural
attributes determine solver input. The exact canonicalization and attributes
are implementation questions, not assumptions about FOGA-11.

For finite-budget analysis, take $\mathcal{X}_B$ to be the set of
phenotypes reachable within budget $B$; assume it is finite. A mutation kernel
$K_t(x'\mid x,s)$ is the conditional distribution of the next phenotype given
parent phenotype $x$ and full pre-mutation state $s$. It includes the selected
structural edit, repair, and canonicalization. Dependence on $t$ or $s$ permits
operator adaptation and history-dependent mutation. The corresponding kernel
on encodings is not necessarily Markov after projecting through $\phi$.

\subsection{Population memory and parent survival}
At proposal epoch $t$, let $P_t$ be the parent population (a multiset of
encodings or phenotypes) and $A_t$ any quality-diversity archive. Let $S_t$
collect scores, lineage information, solver state, and other algorithm
variables that affect the next transition. Thus $M_t=(P_t,A_t,S_t)$ is the
operational memory state. Parent survival is a transition event: for parent
$i$, its survival probability is
\[
 r_{t,i}=\Prb(i\in P_{t+1}\mid M_t,\text{environment at }t).
\]
This is distinct from its parent-selection probability
$\pi_{t,i}=\Prb(i\text{ is selected for a proposal}\mid M_t)$.
Retaining a parent does not imply it is selected, and selection does not
guarantee a useful mutation. The full process may be time-inhomogeneous
Markovian in $(M_t,\theta_t,q_t)$ if that state is sufficient; a niche label
or niche-count vector alone need not be.

\subsection{Hardness and moving target}
Let $\theta_t$ describe the task/environment and $q_t$ the state of the
co-adapting BQ solver. Define true hardness $h_t(x)=h(x;\theta_t,q_t)$,
where a larger value means harder under a prespecified solver-performance
metric. For threshold $\kappa$, the true hard region at epoch $t$ is
\[
 \HARD_t=\{x\in\mathcal{X}_B:h_t(x)\geq\kappa\}.
\]
This definition allows the target to move because either the task or solver
state changes. The event of discovery is evaluated at a declared evaluation
epoch: a proposal $X_t$ hits when $X_t\in\HARD_t$. For the proofs, $\HARD_t$
must be determined before drawing/evaluating $X_t$ (it may be random but
predictable from the pre-proposal history). Whether actual BLADE labels a
candidate ``hard'' by HGB40 alone or by BQ verification is unspecified.

HGB40 supplies an estimate $\widehat h_t(x)$, not the definition of the
latent target. Surrogate ranking can be used by selection, but an estimate of
reachability into $\HARD_t$ requires calibration or a direct verification
model. Without such a condition, true-hardness hit probabilities cannot be
inferred from HGB40 scores.

\subsection{First hit and computational normalization}
Index proposal/evaluation cycles by $i$. Let
\[
 \tau=\inf\{i\geq 1:X_i\in\HARD_i\},\qquad
 C_\tau=\sum_{i=1}^{\tau}c_i,
\]
where $c_i>0$ is the normalized compute cost through evaluating proposal $i$.
The budget-$B$ discovery event is $\{C_\tau\leq B\}$. The cost must include
the work that differs between methods (structural generation, HGB40 inference,
archive operations, and BQ evaluation/verification, as applicable). Counting
generations alone is not computational normalization. The elementary bound
below uses deterministic proposal costs $c_i$; random or method-dependent
costs require a cost-aware hazard analysis and are not silently replaced by
generation counts.

For a policy $\pi$, define its conditional hit hazard
\[
 q_i^\pi=\Prb(X_i\in\HARD_i\mid\mathcal{F}_{i-1}^\pi),
\]
where $\mathcal{F}_{i-1}^\pi$ is the complete pre-proposal history. On a
no-hit history, a lower bound on $q_i^\pi$ is a reachability certificate. A
Markov state representation is one way to calculate this hazard, but the
history-based definition does not require Markovianity.

\subsection{Abstraction versus supplied BLADE-NCO context}
\begin{center}
\begin{tabular}{>{\raggedright\arraybackslash}p{0.22\linewidth}>{\raggedright\arraybackslash}p{0.35\linewidth}>{\raggedright\arraybackslash}p{0.34\linewidth}}
\toprule
Supplied component & Representation here & Unspecified for a BLADE-specific proof \\
\midrule
FOGA-11 mutation & Phenotype kernel $K_t$ & Edit set, operator probabilities, repairs, and transition support \\
Inherited parents; Elite-QD1000 (P=T) & Parent population $P_t$ and archive $A_t$ & Meaning of $P=T$, archive update, parent survival, and selection probabilities \\
HGB40 and BQ co-adaptation & Surrogate $\widehat h_t$ and true $h(x;\theta_t,q_t)$ & Hardness metric/threshold, calibration, solver update timing, and validation cost \\
Fixed generator and HAC & Comparator proposal kernels & Exact proposal laws, gradient access, stopping rules, and per-proposal compute \\
Finite budget & Cumulative cost $C_\tau$ & Which compute resources are charged and how variable costs are normalized \\
\bottomrule
\end{tabular}
\end{center}
Accordingly, the formal model is a deliberately explicit abstraction, not a
claim that BLADE-NCO has exactly these transition probabilities or state
variables. The missing implementation semantics are recorded as open
assumptions in \texttt{assumption\_register.md}.

\subsection{Niche aggregation and the Markov condition}
Suppose $g$ maps the full state to a niche label or niche-count summary. The
aggregated process $g(M_t)$ is Markov only if states with the same aggregate
have equal transition probabilities into every aggregate block (strong
lumpability). Otherwise the aggregate hides parent identity, within-niche
structure, or solver/archive history that predicts the next transition. An
exact failure example is given in Section~\ref{sec:counterexamples}.
\section{Candidate propositions}

\subsection{Candidate 1: finite-window reachability comparison}
\mstatus{CONDITIONAL} \istatus{PARTIAL} \nstatus{NOT CLAIMED} The following
is the strongest surviving candidate.

\begin{proposition}[Memory advantage on a stable target window]
\label{prop:window}
Suppose the true hard region attached to the declared discovery event is
constant, $\HARD_j=\HARD$, for a window of $W$ comparison opportunities and
each opportunity has the same deterministic compute cost
$c>0$. On every no-hit history in that window, suppose a memory policy $M$
has conditional hit hazard at least $p_M$, and a comparator $G$ has
conditional hit hazard at most $p_G$, with $p_M>p_G$. Then, for
$1\leq n\leq W$,
\[
 \Prb(\tau_M\leq n)\geq 1-(1-p_M)^n
 > 1-(1-p_G)^n\geq \Prb(\tau_G\leq n),
\]
and, for $2\leq n\leq W$,
\[
 \Ex[c\min(\tau_M,n)]
 \leq c\sum_{j=0}^{n-1}(1-p_M)^j
 < c\sum_{j=0}^{n-1}(1-p_G)^j
 \leq \Ex[c\min(\tau_G,n)].
\]
For a generic policy that actually samples a parent, one sufficient
derivation of $p_M$ is the following: on each no-hit history,
with conditional probability at least $r$ a useful retained parent is
available; given that event it is selected with probability at least
$\gamma$; conditional on that selection, its mutation/evaluation reaches the
true current hard region with probability at least $\alpha$. Then
$p_M\geq r\gamma\alpha$. No independence among these events is assumed;
the bounds are conditional in the stated order.

That random-selection factorization is not the S6-AE schedule. For the
verified deterministic Cartesian batch, every retained parent receives all
11 operators. If $r$ lower-bounds availability of a useful parent and
$\alpha_{\rm batch}$ lower-bounds, under the explicitly declared RNG
ensemble, the conditional probability that its scheduled operator batch
generates the declared hit, then the appropriate sufficient bound is
$p_M\geq r\alpha_{\rm batch}$. With roots and all seeds conditioned to their
frozen values, this is a statement about a realized batch, not an estimated
random parent/operator-selection probability.
\end{proposition}

\paragraph{Importance.} It turns a structural-memory claim into a finite-budget,
equal-cost comparison that explicitly accounts for parent access and mutation
reachability. It also makes clear that survival alone is not a sufficient
mechanism.

\paragraph{Assumptions and realism.} A piecewise-stable target window and
equal-cost opportunities are useful first approximations. In S6-AE the BQ
checkpoint is fixed during mutation/acquisition in cycle $k$ and changes only
after $T_k$ is frozen and labeled, but comparison windows spanning adaptations
can still move $\HARD_k$. HGB40, archive, exact-label, and BQ-training costs
also differ from a single mutation attempt. The deterministic 10,000-by-11
schedule and $P=T$ update are verified; the history-uniform
$r\alpha_{\rm batch}$ bound, scientific threshold, and matched comparator
cost/hazard remain \istatus{OPEN}.

\paragraph{Proof approach.} Induction on the no-hit survival probability,
then the identity expressing expected truncated hitting time as the sum of
survival probabilities. It requires no conditional independence and no
Markov approximation. Proof details are in \texttt{proof\_attempts.tex}.

\paragraph{Counterexample and HAC distinction.} If memory's conditional
hazard is lower because its retained parents are stale, or if a gradient
proposal has larger hazard per compute unit, the premise fails. The exact
examples in \texttt{counterexamples.tex} show both. Thus this is a conditional
criterion that can distinguish BLADE from fixed generation or HAC only after
their hazards and costs are measured; it does not establish that BLADE wins.

\subsection{Candidate 2: inherited memory never hurts}
\mstatus{DISPROVED} \istatus{NOT APPLICABLE} \nstatus{NOT CLAIMED} The
proposed blanket claim is that retaining historical
parents cannot increase first-hitting cost relative to a fixed generator.
This is false even with two structural states and one target change. The
stale-memory example in Section~\ref{sec:counterexamples} has expected cost
$100$ for inherited memory versus $2$ for a fixed uniform generator. The
failure is not a technical gap in Candidate 1: the required hazard dominance
is reversed.

\paragraph{Importance and proof approach.} A universal no-harm theorem would
have supported the research objective directly, but the exact counterexample
rules it out. No proof is possible without additional target-drift,
forgetting, or reachability assumptions. These extra assumptions would need
to be stated and independently justified; they cannot be inferred from the
word ``memory.''

\paragraph{Assumptions and realism.} The universal claim would require memory never to reduce true-hardness reachability per unit cost for any target shift or parent-selection history. That is not guaranteed by inherited populations and is contradicted by the stale-memory example; realism for BLADE-NCO is low absent explicit target-aware forgetting and verified hazard bounds.

\paragraph{Strongest proof approach.} An exact finite-state counterexample proves the universal claim false. A positive restricted result would need additional drift and reachability premises, which are instead stated explicitly in Candidate 1.

\paragraph{HAC distinction.} Nothing in this false claim compares BLADE to
HAC. HAC can outperform a poorly aligned retained population, as the path
example in Section~\ref{sec:counterexamples} demonstrates.

\subsection{Candidate 3: diversity always improves discovery}
\mstatus{DISPROVED} \istatus{NOT APPLICABLE} \nstatus{NOT CLAIMED} The
proposed blanket claim is that preserving a more
diverse population always reduces first-hitting time. In a three-state model,
a diverse parent selector spends half of each equal-cost proposal on an
irrelevant niche and has expected hitting cost $2$, while a focused selector
hits in cost $1$. The exact construction is in Section~\ref{sec:counterexamples}.

\paragraph{Importance and assumptions.} Diversity can preserve useful future
routes, but its benefit depends on niche relevance, selection allocation,
target drift, and cost. A monotonicity proof from diversity cardinality alone
cannot work. This counterclaim concerns a universal guarantee, not a claim
that Elite-QD1000 behaves like the toy selector.

\paragraph{Realism and strongest proof approach.} No selection-allocation or niche-relevance condition is supplied for Elite-QD1000, so the premise is not verified for BLADE-NCO. The exact three-state construction is a complete disproof of the universal claim; a positive diversity theorem would require explicit target-access conditions. This does not rank BLADE against HAC.

\paragraph{HAC distinction.} It provides no general ranking against HAC.
With equal compute, a focused or gradient-guided policy may spend more
proposals on the current reachable hard region; the actual ranking depends
on transition probabilities and costs, neither of which is specified here.

\section{Proof attempt for the strongest candidate}

\subsection{History-based survival lemma}
\mstatus{PROVED} \istatus{PARTIAL} \nstatus{NOT CLAIMED} The lemma below is
a standard finite-horizon hitting bound,
included as the proof engine rather than as a novelty claim.

Let $Y_j$ be the indicator of the declared discovery event at comparison
opportunity $j$ and let $\tau=\inf\{j\geq1:Y_j=1\}$. The event
$\{\tau\geq j\}$ is determined by the history before opportunity $j$.
Assume there are deterministic numbers $p_j\in[0,1]$ such that, on every
pre-opportunity history with no previous hit,
\[
 \Prb(Y_j=1\mid\mathcal{F}_{j-1})\geq p_j.
\]
This is a history-uniform conditional hazard condition. It allows the target
attached to the event to vary with $j$. The index $j$ is deliberately not the
BQ cycle index $k$: it may represent a serialized mutation attempt, a full
generation event $D_k^{\rm gen}$, or an acquisition event
$D_k^{\rm acq}$, provided the choice and filtration are fixed.

\begin{lemma}[Non-stationary survival and cost bound]
For every $n\geq1$,
\[
 \Prb(\tau>n)\leq\prod_{i=1}^{n}(1-p_i).
\]
If opportunity costs $c_i$ are deterministic, then
\[
 \Ex[C_\tau]
 \leq \sum_{i\geq1}c_i\prod_{j=1}^{i-1}(1-p_j),
\]
with the inequality valid also when the right side is infinite. For a budget $B$, assume
$\sum_{i=1}^\infty c_i=\infty$ and let
$n(B)=\max\{n\geq0: \sum_{i=1}^{n}c_i\leq B\}$, with empty sum zero.
\[
 \Prb(C_\tau\leq B)\geq
 1-\prod_{i=1}^{n(B)}(1-p_i).
\]
\end{lemma}

\begin{proof}
Write $S_i=\{\tau\geq i\}$ for survival to the start of opportunity $i$.
Since $S_i\in\mathcal{F}_{i-1}$,
\begin{align*}
 \Prb(S_{i+1})
 &=\Ex\!\left[\mathbf{1}_{S_i}
       \Prb(Y_i=0\mid\mathcal{F}_{i-1})\right] \\
 &\leq (1-p_i)\Prb(S_i).
\end{align*}
Starting from $\Prb(S_1)=1$ and iterating gives
$\Prb(S_{n+1})\leq\prod_{i=1}^{n}(1-p_i)$, which is the first claim.
For the cost, use the nonnegative-sum identity
\[
 C_\tau=\sum_{i\geq1}c_i\mathbf{1}_{\{\tau\geq i\}},
 \qquad
 \Ex[C_\tau]=\sum_{i\geq1}c_i\Prb(S_i),
\]
and apply the survival bound termwise. Finally, completion of any of the
first $n(B)$ opportunities costs at most $B$, so budget success is at least the
probability of a hit among those opportunities. The first claim supplies the
stated lower bound.
\end{proof}

\subsection{Deriving the window comparison}
For Candidate~1, apply the survival lemma on the $W$-opportunity stable-target
window. The memory policy satisfies $p_j=p_M$ and the comparator satisfies
an upper hazard bound $q_j^G\leq p_G$ on every no-hit history. The same
induction, with the inequality reversed for $G$, yields
\[
 \Prb(\tau_M>j)\leq(1-p_M)^j,
 \qquad
 \Prb(\tau_G>j)\geq(1-p_G)^j.
\]
Therefore, for $n\leq W$,
\[
 \Prb(\tau_M\leq n)\geq1-(1-p_M)^n
 >1-(1-p_G)^n\geq\Prb(\tau_G\leq n).
\]
For $2\leq n\leq W$, using
$\Ex[\min(\tau,n)]=\sum_{j=0}^{n-1}\Prb(\tau>j)$ and multiplying by the
common opportunity cost $c$ gives the strict truncated compute-cost comparison in
Proposition~\ref{prop:window}. At $n=1$ both policies pay the first opportunity
cost in the truncated-cost convention, so only the success-probability
comparison is strict there.

\subsection{Dependency audit}
\begin{itemize}
\item \textbf{Conditional independence:} not used. The proof requires
history-conditional lower/upper hazard bounds; dependence across offspring or
generations is allowed if those bounds remain true after every no-hit history.
\item \textbf{Markovianity:} not used in the hitting proof. It would only be
needed if one chooses a finite-state transition matrix to calculate the
hazards. A niche aggregation is valid for that purpose only under lumpability.
\item \textbf{Mutation bounds:} the quantity $\alpha$ must be a lower bound
on reaching the true current hard region conditional on the useful parent
being selected. A graph path's existence alone is insufficient; its
probability and compute cost matter.
\item \textbf{Parent retention and scheduling:} a lower bound $r$ must concern
availability of a genuinely useful parent at the declared opportunity.
Historical presence before a target shift does not establish continued
usefulness afterward. The S6-AE parent/operator schedule is deterministic:
conditional on $P_k$, each parent receives all 11 operators. Hence no random
parent-selection factor $\gamma$ may be inserted; a batch reachability bound
$\alpha_{\rm batch}$ must be derived under the declared RNG ensemble.
\item \textbf{Surrogate accuracy:} no accuracy is needed for the abstract
survival lemma. To pass from $D_k^{\rm gen}$ to
$D_k^{\rm acq}$, a retention condition connecting $\widehat h_k$ and the QD
rule to true $h_k^{\rm BQ}$ and future reachability is essential. Perfect
current-hardness prediction alone is insufficient when tied candidates have
different mutation futures. Without such a bridge, the implementation-level
step is OPEN.
\item \textbf{Cost normalization:} the comparison assumes equal deterministic
cost $c$ per opportunity. If BLADE archive updates, HGB40 inference, or BQ
validation change cost, opportunity-count dominance does not imply compute-cost
dominance.
\end{itemize}

\subsection{Failure boundary}
\mstatus{OPEN for a BLADE-specific dominance claim}
\istatus{PARTIAL} \nstatus{OPEN} The proof does not establish that BLADE's conditional hazard is
above the fixed-generator or HAC hazard. It fails to produce a BLADE result
at the exact step where one would need verified bounds on $r$ and
$\alpha_{\rm batch}$ (or a direct generation/acquisition hazard bound) and a
comparator upper bound $p_G$, all in the same compute units. The
counterexamples show these inequalities can reverse. No unproved transition
from a structural-memory narrative to hazard dominance is made.

\section{Adversarial finite-state checks}
The examples below are exact toy models. They test universal statements and
assumptions; none is asserted to be the BLADE-NCO implementation.

\subsection{Stale memory after a target shift}
\label{sec:counterexamples}
\status{DISPROVED} for the claim that historical memory cannot hurt.
Take $\mathcal{X}=\{a,b\}$. Before a change the hard region is
$\HARD_0=\{a\}$; just before the first measured proposal it shifts to
$\HARD_i=\{b\}$ for all $i\geq1$. The inherited population contains only
$a$. The memory policy mutates from $a$ to $a$ with probability $0.99$ and to
$b$ with probability $0.01$, retaining $a$ after a miss. Each proposal costs
one unit. Its hitting time is geometric with success probability $0.01$,
so $\Ex[\tau_M]=100$ and
\[
 \Prb(\tau_M\leq B)=1-0.99^B.
\]
A fixed generator that independently samples $a$ or $b$ uniformly every
proposal has $\Ex[\tau_G]=2$ and success probability $1-0.5^B$. Thus the
memory policy is strictly worse at every positive integer budget. The target
has undergone a single non-stationary shift; the failure is stale structural
support, not an independence artifact.

\subsection{Gradient adaptation can beat structural mutation}
\status{DISPROVED} for any universal claim that structural mutation
outperforms HAC-style gradient adaptation. Let
$\mathcal{X}=\{0,1,\ldots,L\}$ and $\HARD=\{L\}$. A structural mutation
from $i<L$ advances to $i+1$ with probability $\varepsilon$ and otherwise
stays at $i$. Starting at $0$, its expected number of unit-cost proposals to
hit $L$ is $L/\varepsilon$. A gradient-informed comparator that advances
deterministically one step per unit-cost proposal hits in $L$ proposals.
For any $0<\varepsilon<1$, the structural policy is slower. This is a toy
landscape with an informative gradient, not a model of the actual HAC
implementation; it proves that operator labels alone imply no ranking.

\subsection{Diversity can allocate effort away from the target}
\status{DISPROVED} for the claim that increased diversity always improves
first hitting time. Let $\mathcal{X}=\{s,u,g\}$, target $\HARD=\{g\}$,
and one proposal cost one unit. The diverse parent set contains distinct
niches $s$ and $u$; parent selection is uniform. Mutation from $s$ reaches
$g$ surely, whereas mutation from $u$ stays at $u$. Hence the per-proposal
hit probability is $1/2$ and expected hitting cost is $2$. A focused policy
that selects $s$ every time hits in cost $1$. Preserving the irrelevant
niche does not help this target and reduces access to the useful parent.

\subsection{Surrogate error can remove a useful bridge}
\status{DISPROVED} for the claim that surrogate-based elite retention
preserves true-hardness reachability. Consider states $\{a,b,g\}$ and
current target $\HARD=\{g\}$; neither $a$ nor $b$ is itself a hit. From $a$,
mutation reaches $g$ surely; from $b$, mutation stays at $b$ surely. A
one-slot score-elite update retains whichever of $a,b$ has the larger
surrogate score. If the true scores of $a,b$ are equal but HGB40
systematically overestimates $b$, the update retains $b$, discards $a$, and
makes the target unreachable. An accurate or reachability-aware selection
would retain $a$. This is a counterexample to a universal error-robustness
claim, not a claim about the precise Elite-QD1000 update rule.

\subsection{Niche aggregation can destroy the Markov property}
\status{DISPROVED} for the assertion that a niche label alone is a Markov
state. Let full states $x_1,x_2$ both map to niche $A$, while $z,w$ map to
distinct niches $Z,W$. From $x_1$, the next state is $z$ with probability
$1$; from $x_2$, the next state is $w$ with probability $1$. Both full states
have the same current aggregate $A$, but the next aggregate is respectively
$Z$ or $W$. The transition probability out of aggregate $A$ is therefore
not determined by the aggregate alone (strong lumpability fails). A Markov
proof using only niche counts would be invalid unless the update rule supplies
the needed lumpability condition or the state is augmented with the hidden
structural/solver information.

\section{Research assessment}
\paragraph{Strongest candidate.}
The useful candidate is the finite-window reachability comparison in Proposition~\ref{prop:window}. It gives an auditable sufficient condition for inherited memory to improve finite-budget discovery on a window where the hard region is stable: the retained lineage must be available, selected, and able to mutate into the \emph{true} current hard region often enough that its per-cost conditional hit hazard exceeds the comparator's. The proof is a survival-probability induction and needs neither independent proposals nor a Markov approximation.

\paragraph{Limitations and novelty.}
The probability argument is standard and is not itself a novel theorem. Its scientific content would have to come from deriving and validating useful lower bounds on FOGA-11 reachability and parent survival, then comparing them with the actual fixed-generator and HAC hazards at equal compute cost. Stage 1.5 verifies the deterministic operator schedule and population update, but the operator-kernel laws under a scientific ensemble, hardness threshold, score-to-retention bridge, matched costs, and HAC transition law remain unspecified. Thus the decisive assumptions cannot yet be checked, and the result does not assert a BLADE advantage.

Stage 1.5 verified the S6-AE batch schedule and population update rule, but not
the required probabilistic hazard bounds, a scientific hard-region threshold,
or a matched HAC transition/cost law. Therefore the mathematical comparison
remains conditional, its BLADE applicability remains partial, and its novelty
remains open.

\paragraph{Testability.}
The result could be tested from existing proposal- and generation-level
records if they identify parent, operator, accepted child, HGB40 score,
acquisition membership, current BQ hardness where evaluated, solver
checkpoint, survival, and comparable cost. The inspected S6-AE artifacts do
record much of this mechanism, but they do not by themselves establish the
history-uniform true-hardness hazard bounds used by the proposition.
\par\istatus{OPEN}\par
No new evaluations, training, or evolutionary runs were performed for this
note.

\paragraph{Unresolved issues.}
The decisive open issue is whether memory yields greater \emph{true-hardness reachability per unit cost} than the fixed generator or HAC after accounting for stale parents, surrogate ranking error, solver co-adaptation, and archive overhead. The exact BLADE transition model and the appropriate cost accounting must be supplied before a BLADE-specific theorem can be attempted.

\section{Stopping point}
Stage 1 formalization, the Stage 1.5 corrections, candidate claims, original
proof attempt, and adversarial checks are recorded. No Stage 2 theorem
development, experimental design, algorithm change, Oracle/BQ evaluation, or
new run is included.

\begin{thebibliography}{9}
\bibitem{Dang2017}
D.-C. Dang, T. Jansen, and P. K. Lehre,
``Populations Can Be Essential in Tracking Dynamic Optima,''
\emph{Algorithmica}, vol. 78, no. 2, pp. 660--680, 2017.
doi: 10.1007/s00453-016-0187-y.

\bibitem{OlivetoZarges2015}
P. S. Oliveto and C. Zarges,
``Analysis of diversity mechanisms for optimisation in dynamic environments with low frequencies of change,''
\emph{Theoretical Computer Science}, vol. 561, part A, pp. 37--56, 2015.
doi: 10.1016/j.tcs.2014.10.028.
\end{thebibliography}

\end{document}
